%% ── Section 3: Methodology ──────────────────────────────────
\section{Methodology}
\label{sec:method}

\subsection{Problem Formulation}

$K = 5$ hospital clients each hold a private dataset
$\mathcal{D}_k = \{(x_i, y_i)\}_{i=1}^{n_k}$ of fundus images $x_i$ with ordinal
grades $y_i \in \{0, 1, 2, 3, 4\}$. A shared EfficientNet-B3 model with weights $w$
is trained over $T$ communication rounds, and no $\mathcal{D}_k$ ever leaves its
client. Because the local grade distributions $P_k(y)$ differ strongly from the
global distribution $P(y)$, the local models drift apart during local training.

\subsection{Limitations of Nominal Objectives}

FL clients are usually trained with cross-entropy (CE). With softmax probabilities
$\mathbf{p}(x) \in \mathbb{R}^5$, the CE loss for true grade $y$ is $-\log p_y(x)$.
This loss does not distinguish a prediction of $y+1$ from one of $y+4$: both are
simply wrong. On a heavily skewed client, for example a clinic with only Grade~0
patients, CE quickly pushes the local model to predict Grade~0 for every image,
and the model loses what it had learned about severe disease.

\subsection{Distance-Aware OrdinalLoss}

To limit this degradation, a hybrid loss, OrdinalLoss, is used during local
training:
\begin{equation}
\mathcal{L}_{\text{ord}}(f(x), y) =
  (1 - \lambda)\,\mathcal{L}_{\text{CE}}(f(x), y)
  + \lambda\,(\hat{y} - y)^2
\label{eq:ordloss}
\end{equation}
where $\mathcal{L}_{\text{CE}}$ is cross-entropy with fixed class weights
$\mathbf{c} = [1.0, 5.0, 1.5, 10.0, 4.0]$ to counter class imbalance, and
$\lambda = 0.5$ balances the two terms. $\hat{y}$ is the expected grade under the
softmax distribution:
\begin{equation}
\hat{y} = \sum_{c=0}^{4} c \cdot p_c(x)
\label{eq:expgrade}
\end{equation}
The penalty $(\hat{y} - y)^2$ grows with the square of the grade error: for a
confident prediction, calling a Grade~4 case Grade~0 costs $16\times$ as much as
calling it Grade~3. Local training is therefore discouraged from predictions far
from the true grade, even when a client's data are dominated by a single grade.

\subsection{Server Validation Set $\mathcal{V}$ and QWK Estimation}

To judge client models without access to their data, the server keeps a small
validation set $\mathcal{V}$ of exactly 50 images: 10 per grade, sampled at random
(seed~42) from the DDR test split, which is separate from all client data. Because
$\mathcal{V}$ is balanced, every grade counts equally in its score. In round $t$,
after receiving the locally updated weights $\tilde{w}_k^t$, the server computes
\begin{equation}
\kappa_k^t =
\mathrm{QWK}\!\left(\{y_i\}_{x_i \in \mathcal{V}},\,
\{\hat{y}_k^t(x_i)\}_{x_i \in \mathcal{V}}\right)
\label{eq:serverqwk}
\end{equation}

\subsection{Adaptive Cosine Schedule $\beta(t)$}

In early rounds, the client models' predictions on $\mathcal{V}$ vary widely, and
each $\kappa_k^t$ is estimated from only 50 images. The aggregation weights are
therefore moved gradually from image counts towards quality scores by a
half-cosine schedule:
\begin{equation}
\beta(t) = \beta_0 \cos\!\left(\frac{\pi (t - 1)}{2 (T - 1)}\right)
\label{eq:beta}
\end{equation}
where $\beta_0 = 0.5$ and $T = 10$. In the first round, image counts and quality
scores count equally ($\beta(1) = 0.5$); in the last round, $\beta(T) = 0$ and only
the quality scores are used.

\subsection{QWK-Weighted Aggregation}

The server-set scores are first floored at $\kappa_{\min}$ and turned into quality
weights with a temperature-scaled softmax:
\begin{equation}
q_k^t = \frac{\exp\!\big(\max(\kappa_k^t, \kappa_{\min}) / \tau\big)}
             {\sum_{j=1}^{K} \exp\!\big(\max(\kappa_j^t, \kappa_{\min}) / \tau\big)}
\label{eq:quality}
\end{equation}
where $\kappa_{\min} = 0.1$ and $\tau = 3$. Each client's aggregation weight blends
this quality weight with the client's share of the images:
\begin{equation}
\omega_k^t = \beta(t)\,\frac{n_k}{\sum_{j=1}^{K} n_j}
           + \big(1 - \beta(t)\big)\, q_k^t
\label{eq:weights}
\end{equation}
and the global model is updated as $w^{t+1} = \sum_{k=1}^{K} \omega_k^t \tilde{w}_k^t$.
A minimum weight of $\omega_{\min} = 0.03$ per client, followed by renormalisation,
is also enforced as a safeguard; with $\beta_0 = 0.5$ and $K = 5$ it never takes
effect, because every client already receives at least
$(1-\beta_0)\,q_{\min} \approx 0.08$. The full procedure is given in
Algorithm~\ref{alg:ordinalfed}.

The scheme is deliberately conservative. No client is ever excluded, because a
single 50-image estimate of $\kappa_k^t$ is too noisy to justify discarding a
hospital's update. With $\tau = 3$ and $\kappa \in [0.1, 1]$, the quality term can
change a client's relative weight by at most a factor of $e^{0.9/3} \approx 1.35$,
so the aggregation is a bounded re-weighting of the image-count weights rather
than a hard filter on client updates.

\begin{algorithm}[t]
\caption{OrdinalFed Training Pipeline}
\label{alg:ordinalfed}
\small
\begin{algorithmic}[1]
\REQUIRE Clients $K$, rounds $T$, local epochs $E$, trade-off $\lambda$,
         initial schedule $\beta_0$, server set $\mathcal{V}$, initial weights $w^0$.
\FOR{$t = 1$ \TO $T$}
  \STATE Broadcast global model $w^t$ to all clients $k \in \{1,\ldots,K\}$
  \FOR{each client $k \in \{1,\ldots,K\}$ \textbf{in parallel}}
    \STATE $\tilde{w}_k^t \leftarrow \textsc{LocalTrain}(w^t, \mathcal{D}_k, E, \lambda)$
           \hfill\textit{(Eq.~\ref{eq:ordloss})}
    \STATE Evaluate validation score $\kappa_k^t$ on $\mathcal{V}$
           \hfill\textit{(Eq.~\ref{eq:serverqwk})}
  \ENDFOR
  \STATE Compute annealing factor $\beta(t)$
         \hfill\textit{(Eq.~\ref{eq:beta})}
  \STATE Compute aggregation weights $\omega_k^t$ for all $k$
         \hfill\textit{(Eqs.~\ref{eq:quality}--\ref{eq:weights})}
  \STATE $w^{t+1} \leftarrow \sum_{k=1}^{K} \omega_k^t \tilde{w}_k^t$
\ENDFOR
\RETURN Final model $w^T$
\end{algorithmic}
\end{algorithm}

% Fig. 2 (QWK trajectories) moved to Sec. V-B, after the paragraph that
% describes it (Reviewer #3, comments 3 and 6).
